% !TEX root = main.tex
% SUPPLEMENTARY FILES (included via % !TEX root = main.tex
% SUPPLEMENTARY FILES (included via % !TEX root = main.tex
% SUPPLEMENTARY FILES (included via % !TEX root = main.tex
% SUPPLEMENTARY FILES (included via \input{supplementary} in main.tex, after
% the References).
%
% Files such as videos, data sets, and software archives can be uploaded to
% ProQuest as "Supplementary Files" along with the dissertation. If you
% upload any, the GSO asks that the dissertation list them in an extra,
% unlabeled appendix that comes after the References. Short sections of
% text, such as computer code, can be included in this appendix as well.
%
% \supplementaryappendix prints the heading "APPENDIX" and the title
% "SUPPLEMENTARY FILES"; give another title in brackets if you prefer:
% \supplementaryappendix[Supplementary Videos and Code].
% List each file with \suppfile{file name}{description}. Type an underscore
% in a file name as \_.
%
% If you upload no supplementary files, delete this file and its \input
% line in main.tex.

\supplementaryappendix

The following files were uploaded to ProQuest with this dissertation as supplementary files. Each description says what the file contains and which part of the dissertation it supports.

\begin{supplementaryfiles}
\suppfile{bisection\_demo.mp4}{A two-minute video of the bisection method (\cref{alg:bisection}) narrowing in on a root.}
\suppfile{wave\_solver.zip}{The MATLAB code used to compute the results in \cref{ch:figtables}, with instructions for running it in the file \texttt{README.txt}.}
\suppfile{symbol\_table.csv}{The data of \cref{table:longsample} as a comma-separated text file.}
\end{supplementaryfiles}
 in main.tex, after
% the References).
%
% Files such as videos, data sets, and software archives can be uploaded to
% ProQuest as "Supplementary Files" along with the dissertation. If you
% upload any, the GSO asks that the dissertation list them in an extra,
% unlabeled appendix that comes after the References. Short sections of
% text, such as computer code, can be included in this appendix as well.
%
% \supplementaryappendix prints the heading "APPENDIX" and the title
% "SUPPLEMENTARY FILES"; give another title in brackets if you prefer:
% \supplementaryappendix[Supplementary Videos and Code].
% List each file with \suppfile{file name}{description}. Type an underscore
% in a file name as \_.
%
% If you upload no supplementary files, delete this file and its \input
% line in main.tex.

\supplementaryappendix

The following files were uploaded to ProQuest with this dissertation as supplementary files. Each description says what the file contains and which part of the dissertation it supports.

\begin{supplementaryfiles}
\suppfile{bisection\_demo.mp4}{A two-minute video of the bisection method (\cref{alg:bisection}) narrowing in on a root.}
\suppfile{wave\_solver.zip}{The MATLAB code used to compute the results in \cref{ch:figtables}, with instructions for running it in the file \texttt{README.txt}.}
\suppfile{symbol\_table.csv}{The data of \cref{table:longsample} as a comma-separated text file.}
\end{supplementaryfiles}
 in main.tex, after
% the References).
%
% Files such as videos, data sets, and software archives can be uploaded to
% ProQuest as "Supplementary Files" along with the dissertation. If you
% upload any, the GSO asks that the dissertation list them in an extra,
% unlabeled appendix that comes after the References. Short sections of
% text, such as computer code, can be included in this appendix as well.
%
% \supplementaryappendix prints the heading "APPENDIX" and the title
% "SUPPLEMENTARY FILES"; give another title in brackets if you prefer:
% \supplementaryappendix[Supplementary Videos and Code].
% List each file with \suppfile{file name}{description}. Type an underscore
% in a file name as \_.
%
% If you upload no supplementary files, delete this file and its \input
% line in main.tex.

\supplementaryappendix

The following files were uploaded to ProQuest with this dissertation as supplementary files. Each description says what the file contains and which part of the dissertation it supports.

\begin{supplementaryfiles}
\suppfile{bisection\_demo.mp4}{A two-minute video of the bisection method (\cref{alg:bisection}) narrowing in on a root.}
\suppfile{wave\_solver.zip}{The MATLAB code used to compute the results in \cref{ch:figtables}, with instructions for running it in the file \texttt{README.txt}.}
\suppfile{symbol\_table.csv}{The data of \cref{table:longsample} as a comma-separated text file.}
\end{supplementaryfiles}
 in main.tex, after
% the References).
%
% Files such as videos, data sets, and software archives can be uploaded to
% ProQuest as "Supplementary Files" along with the dissertation. If you
% upload any, the GSO asks that the dissertation list them in an extra,
% unlabeled appendix that comes after the References. Short sections of
% text, such as computer code, can be included in this appendix as well.
%
% \supplementaryappendix prints the heading "APPENDIX" and the title
% "SUPPLEMENTARY FILES"; give another title in brackets if you prefer:
% \supplementaryappendix[Supplementary Videos and Code].
% List each file with \suppfile{file name}{description}. Type an underscore
% in a file name as \_.
%
% If you upload no supplementary files, delete this file and its \input
% line in main.tex.

\supplementaryappendix

The following files were uploaded to ProQuest with this dissertation as supplementary files. Each description says what the file contains and which part of the dissertation it supports.

\begin{supplementaryfiles}
\suppfile{bisection\_demo.mp4}{A two-minute video of the bisection method (\cref{alg:bisection}) narrowing in on a root.}
\suppfile{wave\_solver.zip}{The MATLAB code used to compute the results in \cref{ch:figtables}, with instructions for running it in the file \texttt{README.txt}.}
\suppfile{symbol\_table.csv}{The data of \cref{table:longsample} as a comma-separated text file.}
\end{supplementaryfiles}
